% 第4回レポート課題（設問2）
% コンパイル例: platex + dvipdfmx（または latexmk -pdfdvi）
\documentclass[a4paper,11pt]{jsarticle}
\usepackage{amsmath,amssymb}
\usepackage{bm}
\usepackage[margin=25mm]{geometry}

\title{第4回レポート課題（設問2）\\
自身の研究に登場する数理モデル：\\
注意状態依存の学習率をもつQ学習モデル}
\author{}
\date{}

\begin{document}
\maketitle

\section{問題設定}

本研究は，\textbf{暗黙的学習（implicit learning）}——参加者が環境の規則性を言語化して報告
できない（気づいていない）にもかかわらず，行動がその規則性に適応していく現象——を対象と
する．具体的には，白色と黒色の2つのランダムドット運動刺激（random dot kinematogram; RDK）
を同時に呈示し，参加者はどちらか一方の色を選んでその運動方向を報告し，報酬（得点）を受け
取る，という課題を用いる．報酬の期待値は「どちらの色が高報酬か」という\textbf{報酬構造の
パターン}（白高配当／黒高配当の2状態）に依存して定められているが，このパターンは参加者に
明示されない．したがって参加者は，試行を重ねる中で報酬フィードバックから報酬構造を
（多くの場合，気づかないまま）学習していくことになる．

このとき問題となるのは，「参加者は\textbf{どれくらいの速さで}報酬構造を学習しているのか」
を，選択行動の系列データから\textbf{定量的に}評価することである．学習の速さを1個のパラメータ
（学習率）として推定できれば，例えば注意が課題から逸れている状態
（マインドワンダリング；本研究では反応時間や報告誤差から各試行に付与する
out-of-the-zone（OOZ）ラベルで操作的に定義する）で学習率が変化するか，といった仮説を
統計的に検定できる．

そのための道具として，強化学習の標準モデルである\textbf{Q学習}を採用する．
ただし本レポートで扱うのは簡単化した\textbf{オラクル版}であり，「参加者は各試行の報酬構造
パターン（白高か黒高か）を正しく知っている」と仮定して，真のパターンをモデルの状態として
直接与える．この仮定は暗黙的学習の前提（パターンは隠されている）とは厳密には整合しないが，
学習率推定の枠組みを最も単純な形で示せるため，比較の基準モデルとして位置づけている．

\section{モデルとアルゴリズム}

\subsection{モデル}

試行 $t$ における状態を $s_t \in \{0,1\}$（$0$: 白高配当，$1$: 黒高配当；オラクル仮定により
真値を与える），行動を $a_t \in \{0,1\}$（$0$: 白を選択，$1$: 黒を選択），報酬を
$r_t \in [0,1]$（得点を正規化したもの）とする．参加者は各状態・行動対の価値
$Q(s,a)$（$2\times2$ の表）を保持し，選択と学習を繰り返すと考える．

\paragraph{選択モデル．}
状態 $s_t$ における黒選択の傾向を
\begin{equation}
  \Delta q_t = \beta\,\bigl(Q(s_t,1) - Q(s_t,0)\bigr) + c
\end{equation}
とおき，黒を選ぶ確率をロジスティック（シグモイド）関数で与える：
\begin{equation}
  P(a_t = 1 \mid s_t) = \sigma(\Delta q_t) = \frac{1}{1+e^{-\Delta q_t}} .
\end{equation}
$\beta > 0$ は逆温度（価値差にどれだけ敏感に選択するか），$c$ は色の選好バイアス
（黒を選びやすい傾向）である．

\paragraph{学習モデル．}
選択後，得られた報酬との\textbf{予測誤差} $\delta_t = r_t - Q(s_t,a_t)$ に比例して
価値を更新する（Rescorla--Wagner 型の更新）：
\begin{equation}
  Q(s_t,a_t) \;\leftarrow\; Q(s_t,a_t) + \alpha_{z_t}\,\bigl(r_t - Q(s_t,a_t)\bigr),
  \qquad
  \alpha_{z_t} =
  \begin{cases}
    \alpha_0 & (z_t = 0:\ \text{非OOZ試行}) \\
    \alpha_1 & (z_t = 1:\ \text{OOZ試行})
  \end{cases}
  \label{eq:update}
\end{equation}
ここで $z_t$ は各試行の注意状態ラベルであり，学習率 $\alpha \in (0,1)$ を注意状態に応じて
切り替える点が本モデルの中心的な仮説である．推定対象のパラメータは
$\bm{\theta} = (\alpha_0, \alpha_1, \beta, c)$ の4個で，
$\log(\alpha_1/\alpha_0)$ が「注意状態が学習の速さに与える効果」の検定統計量となる．

\subsection{パラメータ推定アルゴリズム（MAP推定）}

参加者ごとの選択系列 $a_{1:T}$ に対する対数尤度は，モデルを試行順に前向きに走らせながら
\begin{equation}
  \log L(\bm{\theta})
  = \sum_{t=1}^{T} \log P(a_t \mid s_t;\, \bm{\theta},\, Q_t)
\end{equation}
と逐次的に計算できる（$Q_t$ は式~\eqref{eq:update} で更新される価値表）．
少数試行（1名あたり48試行）での推定を安定させるため，各パラメータに事前分布
\begin{equation}
  \alpha_0, \alpha_1 \sim \mathrm{Beta}(2,2), \qquad
  \beta \sim \mathrm{Gamma}(2,\,3), \qquad
  c \sim \mathcal{N}(0,\,2^2)
\end{equation}
を置き，\textbf{事後確率最大化（MAP: maximum a posteriori）}でパラメータを推定する：
\begin{equation}
  \hat{\bm{\theta}}
  = \arg\max_{\bm{\theta}}\;
    \Bigl[\, \log L(\bm{\theta}) + \log p(\bm{\theta}) \,\Bigr].
\end{equation}
目的関数は非凸で局所解を持ちうるため，パラメータ空間上のグリッド点を初期値とする
\textbf{多点スタート}から準ニュートン法（L-BFGS-B，境界制約付き）による数値最適化を行い，
最良解を採用する．手順をまとめると：
\begin{enumerate}
  \item 初期値グリッド（$\alpha_0, \alpha_1, \beta, c$ の直積）を用意する．
  \item 各初期値から L-BFGS-B で負の対数事後確率を最小化する．
  \item 全初期値の中で最良の解 $\hat{\bm{\theta}}$ を採用する．
  \item 推定された $\hat{\alpha}_0, \hat{\alpha}_1$（および $\log(\hat\alpha_1/\hat\alpha_0)$）を
        参加者間で比較・検定する．
\end{enumerate}

\section{信号処理・パターン認識・機械学習技術との関連に対する考察}

\paragraph{（1）適応フィルタ・信号処理との関連．}
更新式~\eqref{eq:update}は $Q \leftarrow (1-\alpha)Q + \alpha r$ と書き直せる．これは報酬系列
$r_t$ に対する\textbf{指数移動平均}，すなわち一次のIIRローパスフィルタ（漏れ積分器）に
ほかならず，学習率 $\alpha$ はフィルタの忘却の速さ（カットオフ）を決める係数に対応する．
$\alpha$ が大きいほど直近の報酬に追従し（帯域が広く雑音に敏感），小さいほど長期平均に
近づく（平滑だが変化への追従が遅い）．「学習率の推定」とは，参加者の脳が報酬信号に
かけているローパスフィルタの時定数をデータから同定する作業だと言い換えられる．
さらに，式~\eqref{eq:update}は2乗予測誤差 $\delta_t^2$ の確率的勾配降下として導出でき，
これは適応フィルタの \textbf{LMS アルゴリズム}において入力ベクトルが状態のワンホット表現で
ある特殊ケースと一致する．心理学のRescorla--Wagner則と信号処理のLMSが同じ更新式に帰着する
という対応は，分野を越えて同一の数理が現れる好例である．

\paragraph{（2）パターン認識・統計的推定との関連．}
選択モデルは価値差を特徴量とする\textbf{ロジスティック回帰}（2クラス識別のシグモイド
識別器）と同型であり，逆温度 $\beta$ は識別境界の鋭さ，バイアス $c$ は切片に対応する．
またパラメータ推定は，尤度に事前分布を組み合わせるベイズ的な\textbf{MAP推定}であり，
少数データにおける過学習（例えば $\alpha \to 0$ や $\beta \to \infty$ への発散）を
事前分布が正則化項として抑える構造は，パターン認識における正則化付き最尤推定と
まったく同じである．

\paragraph{（3）状態推定・機械学習との関連と，オラクル仮定の限界．}
本モデルはQ学習という\textbf{強化学習}（機械学習の一分野）の枠組みそのものであるが，
より広い視点では状態推定理論との接点を持つ．Rescorla--Wagner型の固定学習率の更新は，
価値（重み）を隠れ状態とみなして追跡する\textbf{カルマンフィルタ}において，事後共分散を
捨ててゲインを定数 $\alpha$ に固定した点推定極限として位置づけられる
（Dayan \& Kakade, 2000; Gershman, 2015）．つまり本モデルは「不確実性の表現を省いた
最小構成の逐次ベイズ推定」と解釈できる．
一方で，本レポートで扱ったオラクル版は「参加者が報酬構造パターンを知っている」という
強い仮定を置いており，パターンが隠されているという暗黙的学習の本質とは緊張関係にある．
この仮定を外すには，観測（ドット運動方向）から隠れた報酬構造を推論する定式化——
部分観測マルコフ決定過程（POMDP）や，方向を連続入力とする関数近似——が必要になり，
これはまさに信号処理・状態推定の技術（ベイズフィルタ，基底関数展開）が行動モデルに
合流する地点である．オラクル版はその出発点となる基準モデルとして意味を持つ．

\end{document}
